\documentclass[11pt,a4paper]{article}
\usepackage[utf8]{inputenc}
\usepackage[T1]{fontenc}
\usepackage[spanish,es-noquoting]{babel}
\usepackage{amsmath,geometry,fancyvrb,xcolor}
\geometry{margin=2.3cm}
\newcommand{\adapt}[1]{\par\medskip\noindent\fcolorbox{orange!70!black}{orange!8}{\parbox{\dimexpr\linewidth-2\fboxsep-2\fboxrule}{\textbf{Adaptación.} #1}}\par\medskip}
\newcommand{\errata}[1]{\par\medskip\noindent\fcolorbox{red!70!black}{red!6}{\parbox{\dimexpr\linewidth-2\fboxsep-2\fboxrule}{\textbf{Errata de la teórica.} #1}}\par\medskip}
\DefineVerbatimEnvironment{code}{Verbatim}{fontsize=\footnotesize,frame=single,framesep=2mm,numbers=left,numbersep=2mm}
\title{Ejercicio 1: explicación del código de los integradores\\\large TP4 -- Dinámica molecular regida por el paso temporal}
\author{Grupo 5}
\date{}
\begin{document}
\maketitle

\section{Contexto común}
El oscilador amortiguado es $m\ddot r=-k r-\gamma\dot r$ con $m=70$ kg, $k=10^4$ N/m, $\gamma=100$ kg/s, $A=1$ m, $t_f=5$ s (\texttt{oscillator.h}). Condiciones iniciales:
\[ r(0)=A,\qquad v(0)=-\frac{A\gamma}{2m}, \]
que son la posición y la derivada en $t=0$ de la solución analítica $r(t)=Ae^{-\gamma t/2m}\cos(\omega t)$.

Cada método es una función \texttt{integrate*(p, dt, steps, emit)} en \texttt{integrators.cpp}. Llama a \texttt{emit(step, r, v)} exactamente \texttt{steps+1} veces (paso 0 con el estado inicial). El motor no calcula errores ni energías: escribe $(t,r,v)$ y el análisis se hace en Python.

La aceleración sale siempre de \texttt{oscForce(p,r,v)/p.m}, o sea $a=(-kr-\gamma v)/m$. Como la fuerza depende de la velocidad, Verlet y Velocity Verlet necesitan despejar una incógnita implícita, y los tres métodos con memoria necesitan un arranque. Esas son las adaptaciones, marcadas en recuadros naranjas.

\section{Euler predictor-corrector (Teórica 4, diap.~23)}
\textbf{Fórmulas de la diapositiva.} Predecir $v^p=v+a\Delta t$, $r^p=r+v\Delta t$; evaluar $a(t+\Delta t)=f(r^p,v^p)/m$; corregir $v=v+a(t+\Delta t)\Delta t$, $r=r+v(t+\Delta t)\Delta t$.

\begin{code}
double r = initialPosition(p);  double v = initialVelocity(p);
emit(0, r, v);
for (long long step = 1; step <= steps; ++step) {
    const double a  = oscForce(p, r, v) / p.m;      // a(t)
    const double vp = v + a * dt;                   // v^p
    const double rp = r + v * dt;                   // r^p
    const double a1 = oscForce(p, rp, vp) / p.m;    // a(t+dt) con predichos
    v += a1 * dt;                                   // v corregida
    r += v * dt;                                    // r usa la v ya corregida
    emit(step, r, v);
}
\end{code}
\textbf{Cómo se lee.} Las líneas 5--6 son el paso 1 (predecir), la 7 el paso 2 (evaluar) y las 8--9 el paso 3 (corregir). Es una traducción directa. \texttt{r += v*dt} va después de actualizar \texttt{v} porque la diapositiva corrige $r$ con $v(t+\Delta t)$.

\textbf{Adaptaciones:} ninguna. Es el esquema de la diap.~23, no el Euler modificado de la diap.~10, cuya posición lleva un término $\Delta t^2 f/2m$.

\section{Verlet original (diap.~13--15)}
\textbf{Fórmulas de la diapositiva.}
\[ r(t+\Delta t)=2r(t)-r(t-\Delta t)+\frac{\Delta t^2}{m}f(t),\qquad v(t)=\frac{r(t+\Delta t)-r(t-\Delta t)}{2\Delta t}. \]

\begin{code}
const double c    = p.gamma * dt / (2.0 * p.m);
const double kdt2 = p.k * dt * dt / p.m;
double r = initialPosition(p);
const double v0 = initialVelocity(p);
const double a0 = oscForce(p, r, v0) / p.m;
double rPrev = r - v0 * dt + 0.5 * a0 * dt * dt;     // r(-dt)
for (long long step = 0; step <= steps; ++step) {
    const double rNext = (2.0*r - rPrev - kdt2*r + c*rPrev) / (1.0 + c);
    emit(step, r, (rNext - rPrev) / (2.0 * dt));
    rPrev = r;
    r = rNext;
}
\end{code}

\adapt{\textbf{(1) Amortiguamiento en forma cerrada (línea 8).} La fuerza es $f=-kr-\gamma v$ y en Verlet $v(t)=(r(t+\Delta t)-r(t-\Delta t))/2\Delta t$ depende de la incógnita $r(t+\Delta t)$. Sustituyendo en la fórmula de la diapositiva:
\[ r_{n+1}=2r_n-r_{n-1}-\frac{k\Delta t^2}{m}r_n-c\,(r_{n+1}-r_{n-1}),\quad c=\frac{\gamma\Delta t}{2m}. \]
Es lineal en $r_{n+1}$, así que se despeja: $r_{n+1}(1+c)=2r_n-r_{n-1}-\frac{k\Delta t^2}{m}r_n+c\,r_{n-1}$. Eso es la línea 8. La diapositiva muestra fuerzas que no dependen de $v$, y este paso no aparece. Hace falta porque este oscilador sí tiene el término $\gamma\dot r$.}

\adapt{\textbf{(2) Arranque (línea 6).} Verlet necesita $r(-\Delta t)$, que no existe. La diap.~14 dice «Euler evaluado en $-\Delta t$» sin dar la fórmula. Se usa $r(-\Delta t)=r-v_0\Delta t+\tfrac12 a_0\Delta t^2$. El término $\tfrac12 a_0\Delta t^2$ es necesario: sin él el primer paso arrastra un error de orden más bajo y el método degrada a orden~1 (lo dice el comentario del código sobre \texttt{integrateVerlet}).}

\adapt{\textbf{(3) Velocidad emitida (línea 9).} La velocidad es la diferencia centrada de la diap.~15 y exige $r(t+\Delta t)$, así que cada iteración calcula un paso de posición de más. Por eso el bucle va de 0 a \texttt{steps} inclusive y se obtiene $v_N$ en el último paso. Verlet no integra $v$, por lo que no hay una $v$ propia del esquema.}

\section{Velocity Verlet (diap.~17)}
\textbf{Fórmulas de la diapositiva.}
\[ r(t+\Delta t)=r+\Delta t\,v+\frac{\Delta t^2}{2m}f(t),\qquad v(t+\Delta t)=v+\frac{\Delta t}{2m}\big(f(t)+f(t+\Delta t)\big), \]
con los pasos intermedios $v(t+\Delta t/2)=v+a(t)\Delta t/2$ y $v(t+\Delta t)=v(t+\Delta t/2)+a(t+\Delta t)\Delta t/2$.

\begin{code}
const double c      = p.gamma * dt / (2.0 * p.m);
const double kdt2m  = p.k * dt / (2.0 * p.m);
double r = initialPosition(p);  double v = initialVelocity(p);
double a = oscForce(p, r, v) / p.m;
emit(0, r, v);
for (long long step = 1; step <= steps; ++step) {
    const double vHalf = v + 0.5 * a * dt;          // paso intermedio 1
    r += vHalf * dt;                                // r(t+dt)
    v = (vHalf - kdt2m * r) / (1.0 + c);            // paso intermedio 2
    a = oscForce(p, r, v) / p.m;                    // a(t+dt), se reutiliza
    emit(step, r, v);
}
\end{code}

\errata{La diapositiva imprime el término de posición como $\Delta t^2 f/m$. El correcto es $\Delta t^2 f/(2m)$. Sale del Taylor de la diap.~15 y es lo que da \texttt{r += vHalf*dt}, porque $\Delta t(v+\tfrac12 a\Delta t)=v\Delta t+\tfrac12 a\Delta t^2$. Con $\Delta t^2 f/m$ el método colapsa a orden~2.}

\adapt{\textbf{Segundo medio paso implícito (línea 9).} El paso intermedio 2 necesita $a(t+\Delta t)=(-kr_{n+1}-\gamma v_{n+1})/m$, que depende de la incógnita $v_{n+1}$. Escribiendo $v_{n+1}=v_{half}+\tfrac{\Delta t}{2}a_{n+1}$ y despejando:
\[ v_{n+1}\Big(1+\frac{\gamma\Delta t}{2m}\Big)=v_{half}-\frac{k\Delta t}{2m}r_{n+1}. \]
La línea 9 es esa ecuación. Se resuelve exacta porque la fuerza es lineal en $v$, y la diapositiva no cubre fuerzas con $v$. Después se recalcula \texttt{a} con la $v$ nueva (línea 10), de modo que el paso siguiente la use como $a(t)$.}

\section{Beeman predictor-corrector (diap.~20)}
\textbf{Fórmulas de la diapositiva.}
\begin{align*}
r(t+\Delta t)&=r+v\Delta t+\tfrac23 a(t)\Delta t^2-\tfrac16 a(t-\Delta t)\Delta t^2,\\
v^{pred}&=v+\tfrac32 a(t)\Delta t-\tfrac12 a(t-\Delta t)\Delta t,\\
v^{corr}&=v+\tfrac13 a(t+\Delta t)\Delta t+\tfrac56 a(t)\Delta t-\tfrac16 a(t-\Delta t)\Delta t,
\end{align*}
con $a(t+\Delta t)$ evaluada con $r(t+\Delta t)$ y $v^{pred}$.

\begin{code}
double a = oscForce(p, r, v) / p.m;
const double rBack = r - v * dt + 0.5 * a * dt * dt;
const double vBack = v - a * dt;
double aPrev = oscForce(p, rBack, vBack) / p.m;      // a(-dt)
emit(0, r, v);
for (long long step = 1; step <= steps; ++step) {
    const double rNew  = r + v*dt + (2.0/3.0)*a*dt*dt - (1.0/6.0)*aPrev*dt*dt;
    const double vPred = v + (3.0/2.0)*a*dt - (1.0/2.0)*aPrev*dt;
    const double aNew  = oscForce(p, rNew, vPred) / p.m;
    const double vNew  = v + (1.0/3.0)*aNew*dt + (5.0/6.0)*a*dt - (1.0/6.0)*aPrev*dt;
    aPrev = a;  a = aNew;  r = rNew;  v = vNew;
    emit(step, r, v);
}
\end{code}
\textbf{Cómo se lee.} Las líneas 7--10 son, una a una, las cuatro ecuaciones de la diapositiva. La línea 11 desplaza la memoria: $a(t)\to a(t-\Delta t)$ y $a(t+\Delta t)\to a(t)$.

\adapt{\textbf{(1) Arranque (líneas 1--4).} Beeman necesita $a(-\Delta t)$. La diapositiva no dice cómo conseguirlo. Se toma un estado Euler hacia atrás ($r(-\Delta t)=r-v\Delta t+\tfrac12 a\Delta t^2$, $v(-\Delta t)=v-a\Delta t$) y se evalúa la fuerza ahí.}

\adapt{\textbf{(2) Una sola evaluación de fuerza por paso (líneas 9 y 11).} $a(t+\Delta t)$ se calcula una vez con $(r_{new},v_{pred})$, como pide la diapositiva, y se guarda para usarla como $a(t)$ en el paso siguiente. No hay una evaluación extra con la $v$ corregida. Por eso Beeman no necesita despejar ninguna ecuación implícita: el predictor de $v$ ya la evita.}

\section{Resumen de adaptaciones}
\begin{center}
\begin{tabular}{@{}l p{4.6cm} p{6.4cm}@{}}
\hline
Método & Adaptación & Motivo \\ \hline
Euler PC & ninguna & Directo de la diap.~23. \\
Verlet & $r_{n+1}$ en forma cerrada & La fuerza depende de $v$, que en Verlet es centrada e implica $r_{n+1}$. \\
Verlet & $r(-\Delta t)$ con $\tfrac12 a_0\Delta t^2$ & Falta el paso previo; sin ese término degrada a orden~1. \\
Verlet & $v$ por diferencia centrada, un paso extra & Es lo que dice la diap.~15; exige $r(t+\Delta t)$. \\
Velocity Verlet & $\Delta t^2/2m$ en vez de $\Delta t^2/m$ & Errata de la diapositiva. \\
Velocity Verlet & $v_{n+1}$ en forma cerrada & $a(t+\Delta t)$ depende de $v(t+\Delta t)$. \\
Beeman & $a(-\Delta t)$ con estado Euler hacia atrás & La diapositiva no da el arranque. \\ \hline
\end{tabular}
\end{center}
\end{document}
